% MAGRED QCOVER2 2026-10-18 12:35 KST
\documentclass{article}
\usepackage{amsmath}
\begin{document}
\section*{Two-sided MCSP source-query obstruction}
Let $N=2^n$, and set $U=(s+1)(n+s)[3(n+s)^2]^s$. At most $U$ truth tables have fan-in-two circuit size at most $s$. A deterministic reduction preserving both MCSP[$s$] decisions must probe at least $N-\log_2 U$ bits even on the zero input: otherwise a NO function agrees with zero on every queried position and produces the identical output. For $s=N^\beta$, $0<\beta<1$ fixed, this forces $N-o(N)$ probes.
A separate randomized sparse-hard argument gives $k=\Omega(N/s)$ source probes for any bounded-error two-sided reduction: a random $t=\Theta(s)$ sparse function has circuit size greater than $s$ with overwhelming probability, while $k$ probes distinguish it from zero with probability at most $kt/N$. This is a source-probe theorem, not a one-tape or RAM running-time lower bound. The independent threshold-descent and model-preservation gaps remain OPEN; P versus NP remains OPEN. Finite integer audit: experiments/meta/2026-10-18_magred_two_sided_query_audit.py.
\nThe sparse count uses $\binom Nt\ge(N/t)^t$, hence $\log_2\binom Nt\ge(8-o(1))sn$, while $\log_2U(n,s)\le(2\beta+o(1))sn$. Conditioning on hard sparse functions removes probability at most $\eta$ and leaves the coupling argument valid.\n\n\section*{Conditional slope-capacity Pareto}\nFor a genuine same-one-tape-model compiler with target length $N'=N^a$, setup exponent $b<1+g$ and target-call cost exponent $a(1+h)$, a source bound $N^{1+g}$ certifies at most $\widehat g=(1+g)/a-1$ target excess. A strict certified slope improvement forces $a\beta'<\beta(1-(a-1)/g)\le\beta$. By the earlier THRES-ENTROPY1 cardinality theorem, an all-YES injective map is then impossible. A surviving one-call compiler must merge YES objects. This is conditional on full one-tape cost and promise preservation.\n\end{document}
% END_OF_VERIFIED_REPORT::MAGRED::QCOVER2::P_VS_NP_OPEN
